\documentclass[11pt,a4paper]{article}
\usepackage[utf8]{inputenc}
\usepackage[T1]{fontenc}
\usepackage{lmodern}
\usepackage{amsmath,amssymb}
\usepackage{booktabs}
\usepackage{geometry}
\usepackage{xr-hyper}
\usepackage[colorlinks=true,linkcolor=blue!50!black,urlcolor=blue!50!black]{hyperref}
\usepackage{xcolor}
\usepackage{tikz}
\usetikzlibrary{arrows.meta,positioning}
\geometry{a4paper,left=2.2cm,right=2.2cm,top=2.2cm,bottom=2.4cm}
\setlength{\parindent}{0pt}
\setlength{\parskip}{0.6ex}
\numberwithin{equation}{section}
\externaldocument[W-]{MG_WALKTHROUGH}

\newcommand{\dg}{^{\dagger}}
\newcommand{\Dpv}{D_{\mathrm{PV}}}
\newcommand{\Cpv}{C_{\mathrm{PV}}}
\newcommand{\Spv}{S_{\mathrm{PV}}}
\newcommand{\Pp}{P_{+}}
\newcommand{\Pm}{P_{-}}
\newcommand{\Gfive}{\Gamma_5}
\newcommand{\AR}{\tilde A_{R}}
\newcommand{\AL}{\tilde A_{L}}
\newcommand{\Pro}{\mathcal P}
\newcommand{\Dc}{\mathcal D_{\mathrm c}}
\newcommand{\Dcpv}{\mathcal D_{\mathrm c}^{\mathrm{PV}}}
\newcommand{\Ce}{C_{e}}
\newcommand{\Co}{C_{o}}
\newcommand{\Cpve}{C_{\mathrm{PV},e}}
\newcommand{\Cpvo}{C_{\mathrm{PV},o}}
\newcommand{\code}[1]{\texttt{\small #1}}

\title{\textsf{\bfseries HMC with the Even-odd Multigrid: for M\"obius Domainwall Fermion}}
\author{CHEN, Wei-Lun}
\date{}

\begin{document}
\maketitle

Note on HMC with the even-odd multigrid solver of RutticeQCD.


This document goes through HMC for M\"obius domain-wall fermions, from the path integral to the
solves of the force, and shows where the multigrid enters. It is the companion of the MG
walkthrough\footnote{\code{docs/mg\_background/MG\_WALKTHROUGH.tex}; references to it are marked
``of the MG walkthrough''.}, which derives the operator, the even-odd forms, the multigrid and the
solve with $D\dg$; the MG note leaves HMC to future work. Nothing here is new with respect to the
code.

\tableofcontents


\section{What we want to compute}
Lattice QCD computes expectation values in the path integral,
\begin{equation}
\langle\mathcal O\rangle=\frac1Z\int\mathcal DU\,\mathcal D\bar\psi\,\mathcal D\psi\;
\mathcal O\,e^{-S_g[U]-\bar\psi D[U]\psi},\qquad
Z=\int\mathcal DU\,\mathcal D\bar\psi\,\mathcal D\psi\;e^{-S_g[U]-\bar\psi D[U]\psi},
\end{equation}
with $U$ the gauge links, $S_g$ the gauge action (here the plaquette action at $\beta=6.0$) and
$\psi$ the quark field. The quark fields are Grassmann numbers, and their integral can be done
exactly:
\begin{equation}
\int\mathcal D\bar\psi\,\mathcal D\psi\;e^{-\bar\psi D\psi}=\det D .
\end{equation}
For two degenerate flavors (up and down) each flavor gives one $\det D$:
\begin{align}
Z&=\int\mathcal DU\;e^{-S_g[U]}\,\det D[U]\,\det D[U]\nonumber\\
&=\int\mathcal DU\;e^{-S_g[U]}\,\det\!\big(D\dg D\big)[U],
\end{align}
where the second line uses that $\det D$ is real, so $(\det D)^2=\det D^*\det D=\det(D\dg D)$.
That $\det D$ is real follows from the exact relation $D_-\Gfive D\dg=D\Gfive D_-\dg$ of
Sec.~\ref{W-sec:g5} of the MG walkthrough, taking determinants ($\det\Gfive=\pm1$ cancels on both sides):
\begin{align}
\det D_-\,\det D\dg&=\det D\,\det D_-\dg\nonumber\\
\det D_-\,(\det D)^*&=\det D\,(\det D_-)^*\nonumber\\
(\det D)^*&=\det D .
\end{align}
The last step divides by $\det D_-$, which is real: since $\Gfive D_-\Gfive=D_-\dg$
(Sec.~\ref{W-sec:g5} of the MG walkthrough),
\begin{equation}
\det D_-=\det(\Gfive D_-\Gfive)=\det D_-\dg=(\det D_-)^* .
\end{equation}

So the problem is to generate gauge fields $U$ with the probability
\begin{equation}
P[U]\;\propto\;e^{-S_g[U]}\det\!\big(D\dg D\big)[U] .
\end{equation}
$\det(D\dg D)$ is the hard part: it couples all sites at once, and computing it directly costs
$O(V^3)$.

\paragraph{The Pauli--Villars ratio.} The domain-wall operator has $L_s$ slices. Only the modes bound
to the walls are the physical quark; the modes in the bulk of the fifth dimension are heavy,
unphysical, and their determinant diverges as $L_s\to\infty$. They are removed by dividing by the
determinant of the same operator at unit mass, the Pauli--Villars operator, whose wall modes are
heavy too:
\begin{equation}
\det\!\big(D\dg D\big)\;\longrightarrow\;\frac{\det\!\big(D\dg D\big)}{\det\!\big(\Dpv\dg\Dpv\big)} .
\end{equation}
This is the ``Ratio'' action.

\section{Pseudofermions: a determinant as a Gaussian integral}
\paragraph{One complex variable.} For a number $a>0$ and a complex variable $\phi=x+iy$,
\begin{align}
\int\frac{dx\,dy}{\pi}\;e^{-\phi^*a^{-1}\phi}
&=\int\frac{dx\,dy}{\pi}\;e^{-(x^2+y^2)/a}\nonumber\\
&=\frac1\pi\Big(\int dx\,e^{-x^2/a}\Big)\Big(\int dy\,e^{-y^2/a}\Big)\nonumber\\
&=\frac1\pi\,\sqrt{\pi a}\,\sqrt{\pi a}=a .
\end{align}
\paragraph{Many variables.} For a Hermitian positive matrix $A=V\Lambda V\dg$ (eigenvalues
$a_k>0$), change variables $\phi=V\chi$; the Jacobian of a unitary change is $1$:
\begin{align}
\int\mathcal D\phi\;e^{-\phi\dg A^{-1}\phi}
&=\int\mathcal D\chi\;e^{-\chi\dg V\dg V\Lambda^{-1}V\dg V\chi}\nonumber\\
&=\int\mathcal D\chi\;e^{-\sum_k|\chi_k|^2/a_k}\nonumber\\
&=\prod_k a_k=\det A .
\end{align}
So a determinant is a Gaussian integral over \emph{bosonic} fields $\phi$, the pseudofermions:
\begin{equation}
\det A=\int\mathcal D\phi\;e^{-\phi\dg A^{-1}\phi} .
\end{equation}
\paragraph{The ratio.} Take $A$ such that $\det A$ is our ratio. With $H=\Gfive D$ (whose
determinant has the modulus of $\det D$),
\begin{align}
A&=H_{\rm PV}\,(H\dg H)^{-1}\,H_{\rm PV}\dg,\nonumber\\
\det A&=\det H_{\rm PV}\,\det(H\dg H)^{-1}\,\det H_{\rm PV}\dg
=\frac{\det(H_{\rm PV}\dg H_{\rm PV})}{\det(H\dg H)}.
\end{align}
Note this is the \emph{inverse} of the ratio; the action uses $A^{-1}$ in the exponent, so
\begin{equation}
\int\mathcal D\phi\;e^{-\phi\dg A\phi}=\det A^{-1}=\frac{\det(H\dg H)}{\det(H_{\rm PV}\dg H_{\rm PV})}
=\frac{\det(D\dg D)}{\det(\Dpv\dg\Dpv)},
\end{equation}
and the pseudofermion action is
\begin{equation}
S_F[U,\phi]=\phi\dg A\phi=\phi\dg H_{\rm PV}(H\dg H)^{-1}H_{\rm PV}\dg\phi .
\end{equation}
The joint distribution of $(U,\phi)$ is $e^{-S_g-S_F}$; integrating out $\phi$ gives back $P[U]$.

\section{Even-odd: only $\mathcal K$ depends on the links}
Factor $D$ into block triangular matrices (LDU):
\begin{equation}
\begin{pmatrix}\Ce&D_{eo}\\D_{oe}&\Co\end{pmatrix}
=\begin{pmatrix}1&D_{eo}\Co^{-1}\\0&1\end{pmatrix}
\begin{pmatrix}S&0\\0&\Co\end{pmatrix}
\begin{pmatrix}1&0\\\Co^{-1}D_{oe}&1\end{pmatrix}.
\end{equation}
Check by multiplying out the last two factors and then the first:
\begin{align}
\begin{pmatrix}S&0\\0&\Co\end{pmatrix}\begin{pmatrix}1&0\\\Co^{-1}D_{oe}&1\end{pmatrix}
&=\begin{pmatrix}S&0\\D_{oe}&\Co\end{pmatrix},\nonumber\\
\begin{pmatrix}1&D_{eo}\Co^{-1}\\0&1\end{pmatrix}\begin{pmatrix}S&0\\D_{oe}&\Co\end{pmatrix}
&=\begin{pmatrix}S+D_{eo}\Co^{-1}D_{oe}&D_{eo}\\D_{oe}&\Co\end{pmatrix}
=\begin{pmatrix}\Ce&D_{eo}\\D_{oe}&\Co\end{pmatrix}.
\end{align}
The triangular factors have determinant $1$, so
\begin{align}
\det D&=\det S\,\det\Co\nonumber\\
&=\det\big(\Ce \mathcal K\big)\det\Co\qquad(\mathcal K=\Ce^{-1}S)\nonumber\\
&=\det\Ce\,\det\Co\,\det \mathcal K .
\end{align}
$\Ce$ and $\Co$ do not depend on the links (Sec.~\ref{W-sec:Cinv} of the MG walkthrough), so they are constants and drop out of
$P[U]$. Only $\det \mathcal K$ is needed, on the even sites alone, half the lattice. The action becomes
\begin{equation}
S_F=\phi_e\dg H_{\rm PV}(H\dg H)^{-1}H_{\rm PV}\dg\phi_e,\qquad H=\Gfive \mathcal K,\ \ H_{\rm PV}=\Gfive \mathcal K_{\rm PV},
\end{equation}
and since $\Gfive$ is Hermitian and unitary, $H\dg H=\mathcal K\dg\Gfive\Gfive \mathcal K=\mathcal K\dg \mathcal K$, so with
$v_1=H_{\rm PV}\dg\phi_e=\mathcal K_{\rm PV}\dg\Gfive\phi_e$,
\begin{equation}
\boxed{\;S_F=v_1\dg\,(\mathcal K\dg \mathcal K)^{-1}\,v_1\;}
\end{equation}
This is the action of the code (Bridge++'s \code{Action\_F\_alt\_Ratio\_eo}).

\section{The idea of HMC}
Sampling $U$ with $P[U]\propto e^{-S[U]}$, $S=S_g+S_F$, by small random changes is hopeless: each
accepted change moves one link a little. HMC moves all links at once along a long path, and still
samples $P[U]$ exactly.

\paragraph{Fictitious momenta.} For every link introduce a momentum $\Pi_\mu(x)$, a traceless
Hermitian $3\times3$ matrix, and the Hamiltonian
\begin{equation}
\mathcal H[U,\Pi]=\tfrac12\sum_{x,\mu}\operatorname{tr}\Pi_\mu(x)^2+S[U] .
\end{equation}
The joint distribution $e^{-\mathcal H}=e^{-\frac12\sum\operatorname{tr}\Pi^2}\,e^{-S[U]}$ is a
product: the momenta are Gaussian and independent of $U$. Sampling $(U,\Pi)$ from $e^{-\mathcal
H}$ and forgetting $\Pi$ samples $U$ from $e^{-S}$.

\paragraph{One trajectory.}
\begin{enumerate}
\item Draw $\Pi$ from its Gaussian and $\phi_e$ by the heat bath (below).
\item Evolve $(U,\Pi)$ for a fictitious time $\tau$ (here $\tau=1$) by Hamilton's equations,
\begin{equation}
\dot U_\mu(x)=i\,\Pi_\mu(x)\,U_\mu(x),\qquad \dot\Pi_\mu(x)=-F_\mu(x),
\end{equation}
with the force $F$ the derivative of $S$ with respect to $U$ (below). Exactly solved, these
conserve $\mathcal H$.
\item The numerical integration does not conserve $\mathcal H$ exactly. Accept the new $U$ with
probability $\min(1,e^{-\Delta\mathcal H})$, $\Delta\mathcal H=\mathcal H_{\rm end}-\mathcal
H_{\rm start}$; otherwise keep the old $U$ (the Metropolis test).
\end{enumerate}
\paragraph{Why the result is exact.} The test makes the method exact for any step size, provided
the integrator has two properties: it is \emph{reversible} (running it with the momenta reversed
brings $(U,-\Pi)$ back) and it \emph{preserves volume} in $(U,\Pi)$ space. Then the proposal
$(U,\Pi)\to(U',\Pi')$ is as likely as its reverse, and the acceptance $\min(1,e^{-\Delta\mathcal H})$
gives detailed balance with $e^{-\mathcal H}$:
\begin{align}
e^{-\mathcal H(U,\Pi)}\min\big(1,e^{-[\mathcal H(U',\Pi')-\mathcal H(U,\Pi)]}\big)
&=\min\big(e^{-\mathcal H(U,\Pi)},e^{-\mathcal H(U',\Pi')}\big)\nonumber\\
&=e^{-\mathcal H(U',\Pi')}\min\big(1,e^{-[\mathcal H(U,\Pi)-\mathcal H(U',\Pi')]}\big),
\end{align}
which is symmetric in the two states. A small step size only raises the acceptance.

\section{The integrator}
\paragraph{Leapfrog.} Split a step of size $\epsilon$ into half a momentum update, a full link
update, and half a momentum update:
\begin{align}
\Pi&\leftarrow\Pi-\tfrac\epsilon2F(U),\nonumber\\
U&\leftarrow e^{i\epsilon\Pi}U,\nonumber\\
\Pi&\leftarrow\Pi-\tfrac\epsilon2F(U).
\end{align}
Each of the three updates changes only $U$ or only $\Pi$ by an amount that depends on the other
one, so each is volume preserving (a shear), and the product of shears is too. Reversing $\Pi$
and running the step again undoes it, update by update. The energy error after a trajectory is
$O(\epsilon^2)$.

\paragraph{Omelyan.} With a parameter $\lambda\approx0.193$,
\begin{equation}
\Pi\leftarrow\Pi-\lambda\epsilon F,\quad U\leftarrow e^{i\frac\epsilon2\Pi}U,\quad
\Pi\leftarrow\Pi-(1-2\lambda)\epsilon F,\quad U\leftarrow e^{i\frac\epsilon2\Pi}U,\quad
\Pi\leftarrow\Pi-\lambda\epsilon F .
\end{equation}
It is still reversible and volume preserving, and $\lambda$ is chosen to make the leading error
term about ten times smaller than that of leapfrog, at two force evaluations per step instead of
one. $e^{i\epsilon\Pi}$ is evaluated by its Taylor series (order 8 in the runs).

\paragraph{Nested steps.} The gauge force is cheap and the fermion force expensive. The runs use
two levels: the fermion force with $n$ Omelyan steps, and within each of its link updates the gauge
force with 2 steps of its own (``\code{number\_of\_steps: [n, 2]}'').

\paragraph{How many solves.} At the fermion level, consecutive momentum updates of neighboring
steps merge into one ($\lambda\epsilon+\lambda\epsilon=2\lambda\epsilon$). Count the force
evaluations of $n$ steps:
\begin{equation}
\underbrace{1}_{\lambda\epsilon\text{ at the start}}+\underbrace{n}_{(1-2\lambda)\epsilon}
+\underbrace{n-1}_{2\lambda\epsilon\text{ between steps}}+\underbrace{1}_{\lambda\epsilon\text{ at the end}}
=2n+1 .
\end{equation}
Each needs one solve with $\mathcal K\dg \mathcal K$, and the energy at the start and the end one more each:
\begin{equation}
\text{solves per trajectory}=2n+3 .
\end{equation}
For $n=20$ this is 43 and for $n=40$ it is 83, the numbers the runs report.

\paragraph{Step size and stability.} The force grows as the quark mass falls (the low modes of $\mathcal K$
come closer to zero). If $\epsilon$ times the largest force frequency exceeds a threshold, the
integrator becomes unstable and $\Delta\mathcal H$ explodes rather than growing smoothly. On the
dynamical $16^4$ configuration at $m=0.01$, $n=10$ gave $\Delta\mathcal H\approx10^3$ (every
trajectory rejected), $n=20$ gave $|\Delta\mathcal H|\approx0.05$ (acceptance $0.98$), and
$n=40$ gave $0.005$.

\section{The heat bath for $\phi_e$}
$\phi_e$ must be drawn from $e^{-S_F}$ at fixed $U$. Draw $\xi_e$ Gaussian, $P(\xi_e)\propto
e^{-|\xi_e|^2}$, and set
\begin{equation}
\phi_e=H_{\rm PV}\,(H_{\rm PV}\dg H_{\rm PV})^{-1}H\dg\,\xi_e .
\end{equation}
This is $\phi_e=(H_{\rm PV}\dg)^{-1}H\dg\xi_e$, written through the normal equation so that CG
can do the solve. Then $S_F=|\xi_e|^2$; the map $\xi_e\mapsto\phi_e$ is linear and does not depend
on $\phi_e$, so its Jacobian is a constant and $\phi_e$ has the distribution $e^{-S_F}$:
\begin{align}
H_{\rm PV}\dg\phi_e&=H_{\rm PV}\dg H_{\rm PV}(H_{\rm PV}\dg H_{\rm PV})^{-1}H\dg\xi_e=H\dg\xi_e,\nonumber\\
S_F&=(H_{\rm PV}\dg\phi_e)\dg(H\dg H)^{-1}(H_{\rm PV}\dg\phi_e)\nonumber\\
&=(H\dg\xi_e)\dg(H\dg H)^{-1}(H\dg\xi_e)\nonumber\\
&=\xi_e\dg H\,H^{-1}H^{-\dagger}\,H\dg\xi_e\nonumber\\
&=\xi_e\dg\xi_e .
\end{align}
This costs one solve with the Pauli--Villars operator ($H_{\rm PV}\dg H_{\rm PV}$, well conditioned:
CG converges fast).

\section{The force}
The force is the change of $S_F$ when the links change. Vary $U$; $\phi_e$ is fixed. With
$\chi=\Gfive\phi_e$, $v_1=\mathcal K_{\rm PV}\dg\chi$ and $v_2=(\mathcal K\dg \mathcal K)^{-1}v_1$, and using
$\delta(M^{-1})=-M^{-1}(\delta M)M^{-1}$:
\begin{align}
\delta S_F&=\delta\big[v_1\dg(\mathcal K\dg \mathcal K)^{-1}v_1\big]\nonumber\\
&=(\delta v_1)\dg(\mathcal K\dg \mathcal K)^{-1}v_1+v_1\dg(\mathcal K\dg \mathcal K)^{-1}\delta v_1+v_1\dg\,\delta\big[(\mathcal K\dg \mathcal K)^{-1}\big]\,v_1\nonumber\\
&=(\delta v_1)\dg v_2+v_2\dg\,\delta v_1-v_1\dg(\mathcal K\dg \mathcal K)^{-1}\,\delta(\mathcal K\dg \mathcal K)\,(\mathcal K\dg \mathcal K)^{-1}v_1\nonumber\\
&=2\,\Re\,(\delta v_1)\dg v_2-v_2\dg\,\delta(\mathcal K\dg \mathcal K)\,v_2\nonumber\\
&=2\,\Re\,(\delta \mathcal K_{\rm PV}\dg\chi)\dg v_2-v_2\dg\big(\delta \mathcal K\dg\,\mathcal K+\mathcal K\dg\,\delta \mathcal K\big)v_2\nonumber\\
&=2\,\Re\,\chi\dg\,\delta \mathcal K_{\rm PV}\,v_2-2\,\Re\,(\mathcal Kv_2)\dg\,\delta \mathcal K\,v_2 .
\end{align}
The fourth line uses that $(\mathcal K\dg\mathcal K)^{-1}$ is Hermitian, so
$v_1\dg(\mathcal K\dg\mathcal K)^{-1}=v_2\dg$, and that $a\dg b+b\dg a=2\Re\,a\dg b$. So one force
evaluation needs exactly one solve, $v_2=(\mathcal K\dg \mathcal K)^{-1}v_1$; the rest are
multiplications. $\delta \mathcal K$ comes from the hops in $S$ alone ($\Ce$ does not depend on the links):
for a variation of one link, only the hops through that link contribute, and the two terms above
become a local outer product of $v_2$, $\mathcal Kv_2$ and $\chi$ at the two ends of the link. With stout
smearing the operator sees the smeared links; the force is taken with respect to them and carried
back to the thin links by the chain rule through the smearing steps.

\section{How precisely to solve}
\paragraph{Energy.} $\mathcal H$ at the start and the end enters $\Delta\mathcal H$ directly, and an
error there biases the Metropolis test. These solves are done to $|r|^2/|b|^2<10^{-24}$.
\paragraph{Force.} An error in the force only changes the trajectory, which the test corrects,
\emph{as long as the integrator stays reversible}. A solve from a zero initial guess is a fixed
function of the links, so reversibility holds to its tolerance; $10^{-16}$ is ample. (Starting the
force solves from the previous solutions, \code{number\_of\_initial\_guess}, saves iterations but
breaks reversibility at the level of the tolerance; it is off by default.)

\section{Solving $(\mathcal K\dg \mathcal K)^{-1}$}
\paragraph{With CG.} $\mathcal K\dg \mathcal K$ is Hermitian positive, so CG solves it directly; each iteration costs
two applications of $\mathcal K$. Mixed precision runs CG in single precision with corrections in double.
With Neff's $\alpha$, Eq.~(\ref{W-eq:Ksim}) of the MG walkthrough gives $\mathcal K=Q_e\mathcal K^{(\alpha)}Q_e^{-1}$ and
\begin{align}
\mathcal K\dg \mathcal K&=Q_e^{-1}\mathcal K^{(\alpha)\dagger}Q_e\,Q_e\mathcal K^{(\alpha)}Q_e^{-1}\nonumber\\
&=Q_e^{-1}\,(Q_e\mathcal K^{(\alpha)})\dg(Q_e\mathcal K^{(\alpha)})\,Q_e^{-1}\nonumber\\
&=Q_e^{-1}\,(\mathcal KQ_e)\dg(\mathcal KQ_e)\,Q_e^{-1}\qquad(Q_e\mathcal K^{(\alpha)}=\mathcal KQ_e),\nonumber\\
(\mathcal K\dg \mathcal K)^{-1}&=Q_e\big[(\mathcal KQ_e)\dg(\mathcal KQ_e)\big]^{-1}Q_e ,
\end{align}
so CG runs on the normal equation of $\mathcal KQ_e$, better conditioned, while the action and the force stay
those of $\mathcal K$ (on $24^4$ at $m=0.001$: 22\% fewer iterations).

\paragraph{With the multigrid: two solves.}\label{sec:twosolves} The multigrid solves with $D$ and $D\dg$, not with
$\mathcal K\dg \mathcal K$. Invert the product, with $\mathcal K=\Ce^{-1}S$:
\begin{align}
(\mathcal K\dg \mathcal K)^{-1}&=\mathcal K^{-1}\mathcal K^{-\dagger}\nonumber\\
&=(\Ce^{-1}S)^{-1}\big[(\Ce^{-1}S)\dg\big]^{-1}\nonumber\\
&=S^{-1}\Ce\;\Ce\dg S^{-\dagger} .
\end{align}
So $z=(\mathcal K\dg \mathcal K)^{-1}r$ is two solves:
\begin{align}
w&=\Ce\dg\,S^{-\dagger}r=\Ce\dg\,\pi D^{-\dagger}\iota\,r &&\text{(with $D\dg$, MG walkthrough Sec.~\ref{W-sec:dagger})},\nonumber\\
z&=S^{-1}\Ce\,w=\pi D^{-1}\iota\,\Ce w &&\text{(with $D$)}.
\end{align}
Each solve stops at $10^{-10}$, while the normal equation needs $10^{-16}$ or $10^{-24}$. So the two
solves are repeated on the residual of the normal equation (iterative refinement):
\begin{align}
r^{(k)}&=b-\mathcal K\dg\mathcal Kz^{(k)},\nonumber\\
z^{(k+1)}&=z^{(k)}+\mathcal K^{-1}\mathcal K^{-\dagger}r^{(k)} ,
\end{align}
where $\mathcal K^{-1}\mathcal K^{-\dagger}r^{(k)}$ stands for the two inexact solves above. If
they reduce the residual by a factor $\epsilon$, then $|r^{(k+1)}|\lesssim\epsilon|r^{(k)}|$,
and each sweep gains about ten digits; the code allows at most eleven sweeps. With $\alpha$ the solve with $D$ uses
$D^{-1}=QD^{(\alpha)-1}$ (the left form $\AL^{(\alpha)}$); the solve with $D\dg$ is unchanged, since
$D^{(\alpha)\dagger}=QD\dg$ only rescales its right-hand side.

\paragraph{When the sweeps stall.} If the sweeps stop gaining (Sec.~\ref{sec:stall}), the solve
is repeated from the current $z$ after a bootstrap of the test vectors for the current links,
twice; if it still stalls, CG finishes it from the current $z$, with a warning. A trajectory is
never lost to the multigrid.

\section{What the multigrid keeps along a trajectory}
The links change a little at every step of the integrator, and the setup is expensive, so the
multigrid is not rebuilt for every solve.
\begin{itemize}
\item The test vectors $y$ are kept from solve to solve (on the host between solves) and improved by
a bootstrap, from the old $y$, when the pseudofermions are drawn; the filter runs only for the
first trajectory.
\item $\Pro$, the coarse operators and the top of the smoother's interval are kept while the links
move (the default, \code{rebuild\_iteration\_growth}~$=1.5$). When the GCR iterations per solve
have grown by half, the test vectors are first improved by a bootstrap for the current links and
then $\Pro$ and the coarse operators are made again. Making $\Pro$ from the kept vectors alone,
for links that have moved, gave 5--9 times the iterations of the $\Pro$ kept (dynamical $16^4$,
$m=0.001$).
\item The cycle is made for each solve from these parts and dropped after it.
\end{itemize}
A kept coarse space is the coarse space of slightly different links. It is still a good
preconditioner; GCR on the operator of the current links keeps the solve exact.

\section{\textcolor{red}{New!!!} The solve with $D\dg$: $\Gfive\mathcal M\Gfive$ or through $D$}
For constant $b$ and $c$ the MG note (Sec.~IV.D) gives the exact relation
$D^{-\dagger}=\Gfive B_5D^{-1}B_5^{-1}\Gfive$ with $B_5=b-cY$, which acts within a site. The solve
with $D\dg$ of HMC can therefore be one solve with $D$ on the right-hand side $B_5^{-1}\Gfive r$
instead of GCR with $\Gfive\mathcal M\Gfive$. Both are in the code; \code{RUTTICE\_B5\_DAGGER=1}
selects the solve through $D$, the default is $\Gfive\mathcal M\Gfive$. Same binary (main
\code{092f332}), $m=0.001$, $\alpha=0.6$, GCR, cache on:
\begin{center}
\begin{tabular}{@{}llrrrr@{}}
\toprule
 & & \multicolumn{2}{c}{$\Gfive\mathcal M\Gfive$} & \multicolumn{2}{c}{through $D$}\\
lattice & inner tolerance & iterations & s & iterations & s\\
\midrule
dynamical $16^4$, 5 trajectories, GCR(2) & $10^{-10}$ & 15208 & 754 & 12328 & 646\\
quenched $24^4$ conf.\ v2, GCR(8) & $10^{-10}$ & 274 & 141 & -- & --\\
 & $10^{-8}$ & 332 & 155 & 299 & 147\\
\phantom{quenched $24^4$} conf.\ 502 & $10^{-10}$ & 240 & 130 & 304 & 145\\
 & $10^{-8}$ & 296 & 144 & 262 & 136\\
\phantom{quenched $24^4$} conf.\ 602 & $10^{-10}$, $10^{-8}$ & -- & -- & -- & --\\
\phantom{quenched $24^4$} conf.\ 702 & $10^{-10}$ & -- & -- & 633 & 312\\
 & $10^{-8}$ & -- & -- & 532 & 286\\
\bottomrule
\end{tabular}
\end{center}
``--'': not finished in 20 minutes. On the dynamical $16^4$ lattice the solve through $D$ takes
19\% fewer iterations and 14\% less time, with $\Delta\mathcal H$ the same to three digits and all
five trajectories accepted. The quenched $24^4$ runs are single steps of length $0.01$ from
quenched configurations ($\Delta\mathcal H$ of $10^3$--$10^4$, all rejected), so they test the
solvers, not HMC. There $\Gfive\mathcal M\Gfive$ does not finish four of the eight runs and the
solve through $D$ three; configuration 602 stalls with both, configuration 702 only with
$\Gfive\mathcal M\Gfive$, and configuration v2 at $10^{-10}$ only through $D$, where the second
sweep's solve with $D$ levels off at $|r|^2/|b|^2\approx7\times10^{-9}$, just above the inner
tolerance. Neither route is the better one on these configurations, so both are kept.

\section{\textcolor{red}{New!!!} Why the solve with $D$ stalls}\label{sec:stall}
The solves of the sweeps that stall are solves with $D$, and their right-hand side comes out of a
solve with $D\dg$. That makes it lie where $D$ is small. Write $S=\sum_i\sigma_iu_iv_i\dg$ for
the singular values $\sigma_i$ of the Schur complement and its left and right singular vectors
$u_i$, $v_i$. The two solves of Sec.~\ref{sec:twosolves} give, up to the site-local factor
$\Ce\Ce\dg$,
\begin{align}
S^{-\dagger}r&=\sum_iu_i\,\frac{v_i\dg r}{\sigma_i}, &
S^{-1}S^{-\dagger}r&=(S\dg S)^{-1}r=\sum_iv_i\,\frac{v_i\dg r}{\sigma_i^2}.
\end{align}
The solve with $D$ thus receives a right-hand side whose $i$-th direction carries $1/\sigma_i$
more weight than in $r$, and its solution $1/\sigma_i^2$ more. On quenched $24^4$ at $m=0.001$ the
smallest $\sigma$ is of order $10^{-3}$. A propagator solve has a right-hand side of ordinary
weights; every solve with $D$ in HMC has this one. In the second and later sweeps it is $Se$,
with $e$ the error the sweep before left, which also lies where $S$ is small.

This right-hand side stalls GCR only where the cycle is poor. GCR($k$) builds a polynomial in
$\AR\mathcal M$ and reduces quickly the part of the residual on which $\AR\mathcal M$ is close to
the identity. At the residual $s$ of a stalled solve we measure $|\AR\mathcal Ms|/|s|=6$--$10$ and
$|\langle s,\AR\mathcal Ms\rangle|/(|s|\,|\AR\mathcal Ms|)\approx10^{-4}$, while the Ritz values
of $\AR\mathcal M$ lie between $0.6$ and $1.7$: there the cycle amplifies instead of inverting,
and $\AR\mathcal M$ is far from normal, so two kept directions cannot lower the residual. How
well the cycle inverts $D$ on its lowest modes is a property of the multigrid; this section only
records where HMC meets it.

The table shows the solves with $D$ of $D\dg D$ solved by sweeps from a random $b$, outside HMC
(setup S of the MG note, $m=0.001$, inner tolerance $10^{-10}$, one GCR solve of up to 2000 steps
in the first sweep and 1000 in the second; iterations, or ``stalls'' with the residual
$|r|^2/|b|^2$ of the solve after 1000 steps where it was recorded).
\begin{center}
\begin{tabular}{@{}llrrr@{}}
\toprule
quenched $24^4$ & & first sweep & second sweep & random right-hand side\\
\midrule
GCR(2) & conf.\ v2 & 15 & stalls ($0.83$) & 7\\
 & conf.\ 502 & 13 & stalls ($0.02$) & 6\\
 & conf.\ 602 & stalls & stalls ($0.98$) & 8\\
 & conf.\ 702 & stalls & stalls ($0.99$) & 10\\
GCR(8) & conf.\ v2 & 13 & 16 & \\
 & conf.\ 602 & 16 & 24 & 8\\
 & conf.\ 702 & 19 & 361 & 8\\
\bottomrule
\end{tabular}
\end{center}
The solves with $D\dg$ of the same sweeps converge in 10--15 steps on every configuration: their
right-hand side is $b$ or the residual of the normal equation. On the dynamical $16^4$ lattice none
of 330 solves with $D$ stalls with GCR(2).

The sweeps add the correction whether or not its solve converged, and a stalled correction makes
the residual larger: on configuration v2 the stalled second sweep raises $|r|^2/|b|^2$ of the normal equation from
$2.8\times10^{-10}$ to $1.8\times10^{-8}$, and every later sweep stalls at that level.

A larger Krylov space gets past the stall. GCR(8) converges on all four configurations, on
configuration 702 only after a plateau at $0.57$ from step 10 to beyond step 50. FGMRES($k$) (\code{solver\_type FGMRES})
needs as many vectors as GCR($k$); in one HMC trajectory on configuration v2 FGMRES(3) takes 305
iterations and 148\,s, GCR(3) 9506 iterations and 2498\,s. Not yet in the code: ending a solve once
it stops gaining, and dropping a stalled correction. Increasing the Chebyshev steps of the
smoother from 8 to 16 cures the second sweep on configuration v2 but in HMC only moves the stall
to a smaller residual ($5\times10^{-14}$ instead of $10^{-10}$).

\section{Memory ($24^4$, $m=0.001$, 20 vectors, levels below)}
The memory table and the measurements below were taken with MINRES on the Pauli--Villars factor,
before the regularized solve; they have not been repeated with it.
\begin{center}
\begin{tabular}{@{}lr@{}}
\toprule
part & GB\\
\midrule
HMC itself (links, smeared levels, momenta, pseudofermion, operators) & 3.3\\
multigrid kept between solves ($\Pro$ 4.75) & 6.6\\
fields of the sweeps & 2.6\\
the solve with $D\dg$ on top (GCR, Schur factors, cycle) & 6.4\\
\midrule
peak (live) & 19.1\\
\bottomrule
\end{tabular}
\end{center}
Mixed CG needs about 9\,GB on the same lattice.

\section{What the measurements say so far}
\begin{itemize}
\item On a thermalized dynamical $16^4$ configuration (60 trajectories at $m=0.05$, then 40 at
$m=0.01$ from the quenched one) the multigrid reproduces the trajectories of CG: the plaquette to
$10^{-11}$, $\Delta\mathcal H$ to three digits. This checks that it solves the right equation.
\item There CG needs about 300 iterations per solve at $m=0.01$ and at $m=0.001$ alike: on a lattice
this small the lowest eigenvalues are held up by the finite volume, not by the mass. The multigrid
needs about 55 iterations, but one of them (a whole cycle) costs about 24 CG iterations, so it is
3--4 times slower per trajectory. Tuning gains at most 6\%.
\item On the quenched $24^4$ configuration CG needs about 6400 iterations per solve and the
multigrid wins, but a trajectory of length 1 there blows up ($\Delta\mathcal H\approx4\times10^9$),
so it is no test of HMC. There are no data yet that show the multigrid faster in HMC; that needs a
volume where CG takes more than about 1300 iterations per solve.
\end{itemize}

\section{Map to the code}
\begin{center}
\small
\begin{tabular}{@{}p{0.36\linewidth}p{0.6\linewidth}@{}}
\toprule
step & where\\
\midrule
HMC: two solves, sweeps, retries & \code{src/hmc/dwf\_multigrid.rs};\newline \code{src/backend/cuda\_bridge/hmc\_multigrid.rs}\\
HMC: action, force, CG with $\alpha$, CG fallback & \code{src/hmc/even\_odd.rs}; \code{src/hmc/solve.rs}\\
\bottomrule
\end{tabular}
\end{center}

\end{document}
